\section{AHSME 1963}
        \begin{enumerate}
        	\item (\textit{AHSME 1963}) Which one of the following points is not on the graph of $y=\dfrac{x}{x+1}$?


 $\textbf{(A)}\ (0,0)\qquad \textbf{(B)}\ \left(-\frac{1}{2},-1\right)\qquad \textbf{(C)}\ \left(\frac{1}{2},\frac{1}{3}\right)\qquad \textbf{(D)}\ (-1,1)\qquad \textbf{(E)}\ (-2,2)$


 
 \textbf{Answer: }D



\textbf{Solution}

Looking at the domain of the graph, $x$ cannot be $-1$. Therefore, the answer is $\boxed{\textbf{(D)}}$.

 


	\item (\textit{AHSME 1963}) let $n=x-y^{x-y}$. Find $n$ when $x=2$ and $y=-2$.


 $\textbf{(A)}\ -14 \qquad \textbf{(B)}\ 0 \qquad \textbf{(C)}\ 1 \qquad \textbf{(D)}\ 18 \qquad \textbf{(E)}\ 256$


 
 \textbf{Answer: }A



\textbf{Solution}

Substitute the variables to determine the value of $n$.
 \[n = 2 - (-2)^{2-(-2)}\]
 \[n = 2 - (-2)^4\]
 \[n = 2 - 16\]
 \[n = -14\]The answer is $\boxed{\textbf{(A)}}$.

 


	\item (\textit{AHSME 1963}) If the reciprocal of $x+1$ is $x-1$, then $x$ equals:


 $\textbf{(A)}\ 0\qquad \textbf{(B)}\ 1\qquad \textbf{(C)}\ -1\qquad \textbf{(D)}\ \pm 1\qquad \textbf{(E)}\ \text{none of these}$


 
 \textbf{Answer: }E



\textbf{Solution}

We form the equation $x+1=\frac{1}{x-1}$.  

 Getting rid of the fraction yields: $x^2-1=1$ $\implies$ $x^2=2$ $\implies$ $x=\pm{\sqrt{2}}=\boxed{\text{E}}$

 



\textbf{Solution 2}

We could plug in every value and see if the statement is satisfied.

 When $x = 0$:
 \[x + 1 = 1\]
 \[x - 1 = -1\]

 Invalid.

 When $x = 1$:
 \[x + 1 = 2\]
 \[x - 1 = 0\]

 Invalid.

 When $x = -1$:
 \[x + 1 = 0\]
 \[x - 1 = -2\]

 Invalid.

 When $x = -1$ or $x = 1$:

 We have already tested both these cases above, and we know that neither satisfy the statement.

 Therefore, none of the answer choices above is the value of $x$. Therefore, the answer is $\boxed{\text{E}}$.

 


	\item (\textit{AHSME 1963}) For what value(s) of $k$ does the pair of equations $y=x^2$ and $y=3x+k$ have two identical solutions?


 $\textbf{(A)}\ \frac{4}{9}\qquad \textbf{(B)}\ -\frac{4}{9}\qquad \textbf{(C)}\ \frac{9}{4}\qquad \textbf{(D)}\ -\frac{9}{4}\qquad \textbf{(E)}\ \pm\frac{9}{4}$


 
 \textbf{Answer: }D



\textbf{Solution}

If the \href{https://artofproblemsolving.com/wiki/index.php?title=System\_of\_equations}{system of equations} has two identical solutions, then only one $(x,y)$ pair will satisfy both equations.

 Substitute $y$ in one equation into another equation.
 \[x^2 = 3x + k\]
 \[x^2 - 3x = k\]\href{https://artofproblemsolving.com/wiki/index.php?title=Complete\_the\_square}{Complete the square} to get
 \[x^2 - 3x + \frac{9}{4} = k + \frac{9}{4}\]
 \[(x - \frac{3}{2})^2 = k + \frac{9}{4}\]In order for the equation to have one solution, the right side must be $0$, so $k = -\frac{9}{4}$, which is answer choice $\boxed{\textbf{(D)}}$.

 


 


	\item (\textit{AHSME 1963}) If $x$ and $\log_{10} x$ are real numbers and $\log_{10} x<0$, then:


 $\textbf{(A)}\ x<0 \qquad \textbf{(B)}\ -1<x<1 \qquad \textbf{(C)}\ 0<x\le 1 \\ \textbf{(D)}\ -1<x<0 \qquad \textbf{(E)}\ 0<x<1$


 
 \textbf{Answer: }E



\textbf{Solution}

Let $\log_{10} x = a$, so $10^a = x$.  Note that $10^0 = 1$ and that $10^a$ increases as $a$ increases.  Since $a < 0$, $10^a < 1$.  However, for all $a$, $10^a > 0$, so $0 < x < 1$.  The answer is $\boxed{\textbf{(E)}}$.

 


	\item (\textit{AHSME 1963}) $\triangle BAD$ is right-angled at $B$. On $AD$ there is a point $C$ for which $AC=CD$ and $AB=BC$. The magnitude of $\angle DAB$ is:


 $\textbf{(A)}\ 67\frac{1}{2}^{\circ}\qquad \textbf{(B)}\ 60^{\circ}\qquad \textbf{(C)}\ 45^{\circ}\qquad \textbf{(D)}\ 30^{\circ}\qquad \textbf{(E)}\ 22\frac{1}{2}^{\circ}$


 
 \textbf{Answer: }B



\textbf{Solution}


                    \begin{figure}[H]
                        centering
                        \includegraphics[width=0.5\textwidth]{Solution/AHSME_I_1963_6_Solution_Figure_1.png}
                    \end{figure}
                    Note that because $C$ is the midpoint of $AD$, $C$ is the \href{https://artofproblemsolving.com/wiki/index.php?title=Circumcenter}{circumcenter} of the triangle, so $BC = AC = AB$.  Thus, $\triangle ABC$ is equilateral, so $\angle DAB = 60^{\circ}$, which is answer choice $\boxed{\textbf{(B)}}$.

 


	\item (\textit{AHSME 1963}) Given the four equations:


 $\textbf{(1)}\ 3y-2x=12 \qquad\textbf{(2)}\ -2x-3y=10 \qquad\textbf{(3)}\ 3y+2x=12 \qquad\textbf{(4)}\ 2y+3x=10$


 The pair representing the perpendicular lines is:


 

$\textbf{(A)}\ \text{(1) and (4)}\qquad \textbf{(B)}\ \text{(1) and (3)}\qquad \textbf{(C)}\ \text{(1) and (2)}\qquad \textbf{(D)}\ \text{(2) and (4)}\qquad \textbf{(E)}\ \text{(2) and (3)}$


 
 \textbf{Answer: }A



\textbf{Solution}

Write each equation in slope-intercept from since the slopes are easier to compare.

 Equation $(1)$ in slope-intercept form is $y = \frac{2}{3}x+4$.

 Equation $(2)$ in slope-intercept form is $y = -\frac{2}{3}x - \frac{10}{3}$.

 Equation $(3)$ in slope-intercept form is $y = -\frac{2}{3}x + 4$.

 Equation $(4)$ in slope-intercept form is $y = -\frac{3}{2}x + 5$.

 Remember that if the two lines are perpendicular, then the product of two slopes equals $-1$.  Equations $(1)$ and $(4)$ satisfy the condition, so the answer is $\boxed{\textbf{(A)}}$.

 


	\item (\textit{AHSME 1963}) The smallest positive integer $x$ for which $1260x=N^3$, where $N$ is an integer, is:


 

$\textbf{(A)}\ 1050 \qquad \textbf{(B)}\ 1260 \qquad \textbf{(C)}\ 1260^2 \qquad \textbf{(D)}\ 7350 \qquad \textbf{(E)}\ 44100$


 
 \textbf{Answer: }D



\textbf{Solution}

Factoring $1260$ results in $2^2 \cdot 3^2 \cdot 5 \cdot 7$.  If an integer $N$ is a perfect cube, then the exponents of all the primes in its \href{https://artofproblemsolving.com/wiki/index.php?title=Prime\_factorization}{prime factorization} are \href{https://artofproblemsolving.com/wiki/index.php?title=Multiples}{multiples} of 3.  Thus, the smallest positive integer that can be multiplied by $1260$ to result in a perfect cube is $2 \cdot 3 \cdot 5^2 \cdot 7^2 = 7350$, which is answer choice $\boxed{\textbf{(D)}}$.

 


	\item (\textit{AHSME 1963}) In the expansion of $\left(a-\dfrac{1}{\sqrt{a}}\right)^7$ the coefficient of $a^{-\dfrac{1}{2}}$ is:


 $\textbf{(A)}\ -7 \qquad \textbf{(B)}\ 7 \qquad \textbf{(C)}\ -21 \qquad \textbf{(D)}\ 21 \qquad \textbf{(E)}\ 35$


 
 \textbf{Answer: }C



\textbf{Solution}

By the \href{https://artofproblemsolving.com/wiki/index.php?title=Binomial\_Theorem}{Binomial Theorem}, each term of the expansion is $\binom{7}{n}(a)^{7-n}(\frac{-1}{\sqrt{a}})^n$.

 We want the exponent of $a$ to be $-\frac{1}{2}$, so 
 \[(7-n)-\frac{1}{2}n=-\frac{1}{2}\]
 \[-\frac{3}{2}n = -\frac{15}{2}\]
 \[n = 5\]

 If $n=5$, then the corresponding term is
 \[\binom{7}{5}(a)^{2}(\frac{-1}{\sqrt{a}})^5\]
 \[-21a^{-\frac{1}{2}}\]

 The answer is $\boxed{\textbf{(C)}}$.

 


	\item (\textit{AHSME 1963}) Point $P$ is taken interior to a square with side-length $a$ and such that is it equally distant from two 
consecutive vertices and from the side opposite these vertices. If $d$ represents the common distance, then $d$ equals:


 $\textbf{(A)}\ \frac{3a}{5}\qquad \textbf{(B)}\ \frac{5a}{8}\qquad \textbf{(C)}\ \frac{3a}{8}\qquad \textbf{(D)}\ \frac{a\sqrt{2}}{2}\qquad \textbf{(E)}\ \frac{a}{2}$


 
 \textbf{Answer: }B



\textbf{Solution}


                    \begin{figure}[H]
                        centering
                        \includegraphics[width=0.5\textwidth]{Solution/AHSME_I_1963_10_Solution_Figure_1.png}
                    \end{figure}
                    Draw a diagram and label it as shown.  Because of SSS Congruency, the two bottom triangles are \href{https://artofproblemsolving.com/wiki/index.php?title=Right\_triangle}{right triangles}.  By the \href{https://artofproblemsolving.com/wiki/index.php?title=Pythagorean\_Theorem}{Pythagorean Theorem},
 \[d^2 = (a-d)^2 + (\frac{a}{2})^2\]
 \[d^2 = a^2 - 2ad + d^2 + \frac{a^2}{4}\]
 \[2ad = \frac{5a^2}{4}\]
 \[d = \frac{5a}{8}\]Thus, the answer is $\boxed{\textbf{(B)}}$.

 


	\item (\textit{AHSME 1963}) The arithmetic mean of a set of $50$ numbers is $38$. If two numbers of the set, namely $45$ and $55$, are discarded, 
the arithmetic mean of the remaining set of numbers is:


 $\textbf{(A)}\ 38.5 \qquad \textbf{(B)}\ 37.5 \qquad \textbf{(C)}\ 37 \qquad \textbf{(D)}\ 36.5 \qquad \textbf{(E)}\ 36$


 
 \textbf{Answer: }B



\textbf{Solution}

If the \href{https://artofproblemsolving.com/wiki/index.php?title=Arithmetic\_mean}{arithmetic mean} of a set of $50$ numbers is $38$, then the sum of the $50$ numbers equals $1900$.  Since $45$ and $55$ are being removed, subtract $100$ to get the sum of the remaining $48$ numbers, which is $1800$.  Therefore, the new mean is $37.5$, which is answer choice $\boxed{\textbf{(B)}}$.

 


	\item (\textit{AHSME 1963}) Three vertices of parallelogram $PQRS$ are $P(-3,-2), Q(1,-5), R(9,1)$ with $P$ and $R$ diagonally opposite. 
The sum of the coordinates of vertex $S$ is:


 $\textbf{(A)}\ 13 \qquad \textbf{(B)}\ 12 \qquad \textbf{(C)}\ 11 \qquad \textbf{(D)}\ 10 \qquad \textbf{(E)}\ 9$


 
 \textbf{Answer: }E



\textbf{Solution}


                    \begin{figure}[H]
                        centering
                        \includegraphics[width=0.5\textwidth]{Solution/AHSME_I_1963_12_Solution_Figure_1.png}
                    \end{figure}
                    Graph the three points on the coordinate grid.  Noting that the opposite sides of a \href{https://artofproblemsolving.com/wiki/index.php?title=Parallelogram}{parallelogram} are congruent and parallel, count boxes to find out that point $S$ is on $(5,4)$.  The sum of the x-coordinates and y-coordinates is $9$, so the answer is $\boxed{\textbf{(E)}}$.

 


	\item (\textit{AHSME 1963}) If $2^a+2^b=3^c+3^d$, the number of integers $a,b,c,d$ which can possibly be negative, is, at most:


 $\textbf{(A)}\ 4 \qquad \textbf{(B)}\ 3 \qquad \textbf{(C)}\ 2 \qquad \textbf{(D)}\ 1 \qquad \textbf{(E)}\ 0$


 
 \textbf{Answer: }E



\textbf{Solution}

Assume $c,d \ge 0$, and WLOG, assume $a<0$ and $a \le b$.  This also takes into account when $b$ is negative.  That means
 \[\frac{1}{2^{-a}} + 2^b = 3^c + 3^d\]Multiply both sides by $2^{-a}$ to get
 \[1 + 2^{-a+b} = 2^{-a} (3^c + 3^d)\]Note that both sides are integers.  If $a \ne b$, then the right side is even while the left side is odd, so equality can not happen.  If $a = b$, then $2^{-a} (3^c + 3^d) = 2$, and since  $a<0$, $a = -1$ and $3^c + 3^d = 1$.  No nonnegative value of $c$ and $d$ works, so equality can not happen.  Thus, $a$ and $b$ can not be negative when $c,d \ge 0$.

 Assume $a,b \ge 0$, and WLOG, assume $c < 0$ and $c \le d$.  This also takes into account when $d$ is negative.  That means
 \[2^a + 2^b = \frac{1}{3^{-c}} + 3^d\]Multiply both sides by $3^{-c}$ to get
 \[3^{-c} (2^a + 2^b) = 1 + 3^{d-c}\]That makes both sides integers.  The left side is congruent to $0$ modulo $3$ while the right side is congruent to $1$ or $2$ modulo $3$, so equality can not happen.  Thus, $c$ and $d$ can not be negative when $a,b \ge 0$.

 Assume $a,c < 0$, and WLOG, let $a \le b$ and $c \le d$.  This also takes into account when $b$ or $d$ is negative.  That means
 \[\frac{1}{2^{-a}} + 2^b = \frac{1}{3^{-c}} + 3^d\]Multiply both sides by $2^{-a} \cdot 3^{-c}$ to get
 \[3^{-c} (1 + 2^{b-a}) = 2^{-a} (1 + 3^{d-c})\]That makes both sides integers.  The left side is congruent to $0$ modulo $3$ while the right side is congruent to $1$ or $2$ modulo $3$, so equality can not happen.  Thus, $a$ and $c$ can not be negative.

 Putting all the information together, none of $a,b,c,d$ can be negative, so the answer is $\boxed{\textbf{(E)}}$.

 


	\item (\textit{AHSME 1963}) Given the equations $x^2+kx+6=0$ and $x^2-kx+6=0$. If, when the roots of the equation are suitably listed, 
each root of the second equation is $5$ more than the corresponding root of the first equation, then $k$ equals:


 $\textbf{(A)}\ 5 \qquad \textbf{(B)}\ -5 \qquad \textbf{(C)}\ 7 \qquad \textbf{(D)}\ -7 \qquad \textbf{(E)}\ \text{none of these}$


 
 \textbf{Answer: }A



\textbf{Solution}

Let the two roots of $x^2 + kx + 6$ be $r$ and $s$.  By \href{https://artofproblemsolving.com/wiki/index.php?title=Vieta\%27s\_Formulas}{Vieta's Formulas},
 \[r+s=-k\]Each root of $x^2 - kx + 6$ is five more than each root of the original, so using Vieta's Formula again yields
 \[r+5+s+5 = k\]
 \[r+s+10=k\]Substitute $r+s$ to get
 \[-k + 10 = k\]
 \[10 = 2k\]
 \[5 = k\]The answer is $\boxed{\textbf{(A)}}$.

 


	\item (\textit{AHSME 1963}) A circle is inscribed in an equilateral triangle, and a square is inscribed in the circle. 
The ratio of the area of the triangle to the area of the square is:


 $\textbf{(A)}\ \sqrt{3}:1\qquad \textbf{(B)}\ \sqrt{3}:\sqrt{2}\qquad \textbf{(C)}\ 3\sqrt{3}:2\qquad \textbf{(D)}\ 3:\sqrt{2}\qquad \textbf{(E)}\ 3:2\sqrt{2}$


 
 \textbf{Answer: }C



\textbf{Solution}


                    \begin{figure}[H]
                        centering
                        \includegraphics[width=0.5\textwidth]{Solution/AHSME_I_1963_15_Solution_Figure_1.png}
                    \end{figure}
                    Let the radius of the circle be $r$.  That means the diameter of the circle is $2r$, so the side length of the square is $r\sqrt{2}$.  Therefore, the area of the square is $2r^2$.

 By using 30-60-90 triangles, half of the side length of an \href{https://artofproblemsolving.com/wiki/index.php?title=Equilateral\_triangle}{equilateral triangle} is $r\sqrt{3}$, so each side is $2r\sqrt{3}$ units long.  Thus, the area of the equilateral triangle is $3r^2\sqrt{3}$.

 The ratio of the area of the equilateral triangle to the area of the square is $\frac{3r^2 \sqrt{3}}{2r^2} = \frac{3 \sqrt{3}}{2}$, so the answer is $\boxed{\textbf{(C)}}$.

 


 



\textbf{Solution 2 (triangle centers)}

Refer to the diagram shown in Solution 1. Notice the dotted line that stretches to the bottom of the square/circle. This is the inradius of the equilateral triangle. If we know this length, we could find the side length of the square, and use the equilateral triangle area formula to find the ratio.

 The inradius of any given triangle is
 \[\frac{\text{Area}}{\text{semiperimeter}}.\]

 We could specify this formula for an equilateral triangle (or if you memorized it for some reason):
 \[r = \frac{\frac{\sqrt{3}s^{2}}{4}}{\frac{3s}{2}} = \frac{s\sqrt{3}}{6}.\]

 Denote this value as $r$. Thus the area of the square is
 \[2r^{2} = 2 \cdot \frac{3s^{2}}{36} = 2 \cdot \frac{s^{2}}{12} = \frac{s^{2}}{6}.\]

 Now ratio this to the area of the equilateral triangle,
 \[\frac{\frac{\sqrt{3}s^{2}}{4}}{\frac{s^{2}}{6}} = 6 \cdot \frac{\sqrt{3}}{4} = \frac{6\sqrt{3}}{4} = \frac{3\sqrt{3}}{2}.\]

 
 \[\rightarrow \boxed{\textbf{(C)}}\]

 


	\item (\textit{AHSME 1963}) Three numbers $a,b,c$, none zero, form an arithmetic progression. Increasing $a$ by $1$ or increasing $c$ by $2$ results 
in a geometric progression. Then $b$ equals:


 $\textbf{(A)}\ 16 \qquad \textbf{(B)}\ 14 \qquad \textbf{(C)}\ 12 \qquad \textbf{(D)}\ 10 \qquad \textbf{(E)}\ 8$


 
 \textbf{Answer: }C



\textbf{Solution}

Let $d$ be the common difference of the \href{https://artofproblemsolving.com/wiki/index.php?title=Arithmetic\_sequence}{arithmetic sequence}, so $a = b-d$ and $c = b+d$.

 Since increasing $a$ by $1$ or $c$ by $2$ results in a \href{https://artofproblemsolving.com/wiki/index.php?title=Geometric\_sequence}{geometric sequence},
 \[\frac{b}{b-d+1} = \frac{b+d}{b}\]
 \[\frac{b}{b-d} = \frac{b+d+2}{b}\]Cross-multiply in both equations to get a \href{https://artofproblemsolving.com/wiki/index.php?title=System\_of\_equations}{system of equations}.
 \[b^2 = b^2 - d^2 + b + d\]
 \[b^2 = b^2 - d^2 + 2b - 2d\]Rearranging terms results in
 \[d^2 = b+d\]
 \[d^2 = 2b-2d\]Substitute and solve for $d$.
 \[b+d = 2b-2d\]
 \[d = \frac{b}{3}\]Finally, substitute $d$ and solve for $b$.  Since $b \ne 0$, dividing by $b$ is allowed.
 \[(\frac{b}{3})^2 = b + \frac{b}{3}\]
 \[\frac{b^2}{9} = \frac{4b}{3}\]
 \[\frac{b}{9} = \frac{4}{3}\]
 \[b = 12\]The answer is $\boxed{\textbf{(C)}}$.

 


 


	\item (\textit{AHSME 1963}) The expression $\dfrac{\dfrac{a}{a+y}+\dfrac{y}{a-y}}{\dfrac{y}{a+y}-\dfrac{a}{a-y}}$, $a$ real, $a\neq 0$, has the value $-1$ for:


 $\textbf{(A)}\ \text{all but two real values of }y \qquad \\ \textbf{(B)}\ \text{only two real values of }y \qquad \\ \textbf{(C)}\ \text{all real values of }y\qquad \\ \textbf{(D)}\ \text{only one real value of }y\qquad \\ \textbf{(E)}\ \text{no real values of }y$


 
 \textbf{Answer: }A



\textbf{Solution}

First, note that $y \ne \pm a$ because that would make the denominator $0$.

 Create common denominators in the complex fraction.
 \[\frac{\frac{a^2-ay}{a^2-y^2} + \frac{ay+y^2}{a^2-y^2}}{\frac{ay-y^2}{a^2-y^2} - \frac{a^2+ay}{a^2-y^2}}\]
 \[\frac{\frac{a^2+y^2}{a^2-y^2}}{\frac{-a^2-y^2}{a^2-y^2}}\]Since $a \ne 0$, $a^2 + y^2 > 0$, so no other restrictions need to be made.
 \[\frac{a^2 + y^2}{a^2 - y^2} \cdot \frac{a^2 - y^2}{-a^2 - y^2}\]
 \[\frac{a^2 + y^2}{-(a^2 + y^2)}\]
 \[-1\]The expression (with suitable restrictions) simplifies to $-1$, so the answer is $\boxed{\textbf{(A)}}$.

 


	\item (\textit{AHSME 1963}) Chord $EF$ is the perpendicular bisector of chord $BC$, intersecting it in $M$. 
Between $B$ and $M$ point $U$ is taken, 
and $EU$ extended meets the circle in $A$. Then, for any selection of $U$, as described, $\triangle EUM$ is similar to:


 
                    \begin{figure}[H]
                        centering
                        \includegraphics[width=0.5\textwidth]{Problem/AHSME_I_1963_18_Problem_Figure_1.png}
                    \end{figure}
                    


 

$\textbf{(A)}\ \triangle EFA \qquad \textbf{(B)}\ \triangle EFC \qquad \textbf{(C)}\ \triangle ABM \qquad \textbf{(D)}\ \triangle ABU \qquad \textbf{(E)}\ \triangle FMC$


 
 \textbf{Answer: }A



\textbf{Solution}

Note that $\triangle EUM$ is a \href{https://artofproblemsolving.com/wiki/index.php?title=Right\_triangle}{right triangle} with one of the angles being $\angle FEA$.  This leads to prediction that $\triangle FEA$ is the similar triangle as it shares an angle, and to prove this, we need to show that $FE$ is a diameter.

 
                    \begin{figure}[H]
                        centering
                        \includegraphics[width=0.5\textwidth]{Solution/AHSME_I_1963_18_Solution_Figure_1.png}
                    \end{figure}
                    

 If we let $D$ be on $FE$ so that $FD = DB$, then by SAS Congruency, $\triangle DMB \cong \triangle DMC$, so $FD = DB = DC$.  Since three points define a circle and point $D$ is equidistant from three points, $D$ is the center, so $FE$ is a diameter.  Therefore, $\angle EAF$ is a right angle, and by AA Similarity, we can confirm that $\triangle EFA \sim \triangle EUM$.  The answer is $\boxed{\textbf{(A)}}$.

 


	\item (\textit{AHSME 1963}) In counting $n$ colored balls, some red and some black, it was found that $49$ of the first $50$ counted were red. 
Thereafter, $7$ out of every $8$ counted were red. If, in all, $90$ \% or more of the balls counted were red, the maximum value of $n$ is:


 $\textbf{(A)}\ 225 \qquad \textbf{(B)}\ 210 \qquad \textbf{(C)}\ 200 \qquad \textbf{(D)}\ 180 \qquad \textbf{(E)}\ 175$


 
 \textbf{Answer: }B



\textbf{Solution}

The desired percentage of red balls is more than $90$ percent, so write an \href{https://artofproblemsolving.com/wiki/index.php?title=Inequality}{inequality}.

 
 \[\frac{49+7x}{50+8x} \ge 0.9\]

 Since $x >0$, the sign does not need to be swapped after multiplying both sides by $50+8x$.

 
 \[49+7x \ge 45+7.2x\]
 \[4 \ge 0.2x\]
 \[20 \ge x\]

 Thus, up to $20$ batches of balls can be used, so a total of $20 \cdot 8 + 50 = 210$ balls can be counted while satisfying the requirements.  The answer is $\boxed{\textbf{(B)}}$.

 


	\item (\textit{AHSME 1963}) Two men at points $R$ and $S$, $76$ miles apart, set out at the same time to walk towards each other. 
The man at $R$ walks uniformly at the rate of $4\frac{1}{2}$ miles per hour; the man at $S$ walks at the constant 
rate of $3\frac{1}{4}$ miles per hour for the first hour, at $3\frac{3}{4}$ miles per hour for the second hour,
and so on, in arithmetic progression. If the men meet $x$ miles nearer $R$ than $S$ in an integral number of hours, then $x$ is:


 $\textbf{(A)}\ 10 \qquad \textbf{(B)}\ 8 \qquad \textbf{(C)}\ 6 \qquad \textbf{(D)}\ 4 \qquad \textbf{(E)}\ 2$


 
 \textbf{Answer: }D



\textbf{Solution}

First, find the number of hours it takes for the two to meet together.  After $h$ hours, the person at $R$ walks $4.5h$ miles.  In the same amount of time, the person at $S$ has been walking at $3.25+0.5(h-1)$ mph for the past hour, so the person walks $\frac{h(6.5+0.5(h-1))}{2}$ miles.

 In order for both to meet, the sum of both of the distances walked must total $76$ miles, so
 \[4.5h + \frac{h(6+0.5h)}{2} = 76\]
 \[4.5h + 3h + 0.25h^2 = 76\]
 \[0.25h^2 + 7.5h - 76 = 0\]
 \[h^2 + 30h - 304 = 0\]
 \[(h + 38)(h - 8) = 0\]Since $h$ must be positive, $h = 8$.  Because it takes $8$ hours to meet, the person from $R$ traveled $36$ miles while the person from $S$ traveled $40$ miles.  Thus, they are $4$ miles closer to $R$ than $S$, so the answer is $\boxed{\textbf{(D)}}$.

 


	\item (\textit{AHSME 1963}) The expression $x^2-y^2-z^2+2yz+x+y-z$ has:


 $\textbf{(A)}\ \text{no linear factor with integer coefficients and integer exponents} \qquad \\ \textbf{(B)}\ \text{the factor }-x+y+z \qquad \\ \textbf{(C)}\ \text{the factor }x-y-z+1 \qquad \\ \textbf{(D)}\ \text{the factor }x+y-z+1 \qquad \\ \textbf{(E)}\ \text{the factor }x-y+z+1$


 
 \textbf{Answer: }E



\textbf{Solution}

Factor the perfect square trinomial.
 \[x^2 - (y-z)^2 + x + y - z\]Factor the \href{https://artofproblemsolving.com/wiki/index.php?title=Difference\_of\_squares}{difference of squares}.
 \[(x+y-z)(x-y+z) + x + y - z\]Factor by grouping.
 \[(x+y-z)(x-y+z+1)\]The answer is $\boxed{\textbf{(E)}}$.

 


	\item (\textit{AHSME 1963}) Acute-angled $\triangle ABC$ is inscribed in a circle with center at $O$; $\stackrel \frown {AB} = 120$ and $\stackrel \frown {BC} = 72$. 
$A$ point $E$ is taken in minor arc $AC$ such that $OE$ is perpendicular to $AC$. Then the ratio of the magnitudes of $\angle OBE$ and $\angle BAC$ is:


 $\textbf{(A)}\ \frac{5}{18}\qquad \textbf{(B)}\ \frac{2}{9}\qquad \textbf{(C)}\ \frac{1}{4}\qquad \textbf{(D)}\ \frac{1}{3}\qquad \textbf{(E)}\ \frac{4}{9}$


 
 \textbf{Answer: }D



\textbf{Solution}


                    \begin{figure}[H]
                        centering
                        \includegraphics[width=0.5\textwidth]{Solution/AHSME_I_1963_22_Solution_Figure_1.png}
                    \end{figure}
                    

 Because $\stackrel \frown {AB} = 120^\circ$ and $\stackrel \frown {BC} = 72^\circ$, $\stackrel \frown {AC} = 168^\circ$.  Also, $OA = OC$ and $OE \perp AC$, so $\angle AOE = \angle COE = 84^\circ$.  Since $\angle BOC = 72^\circ$, $\angle BOE = 156^\circ$.  Finally, $\triangle BOE$ is an \href{https://artofproblemsolving.com/wiki/index.php?title=Isosceles\_triangle}{isosceles triangle}, so $\angle OBE = 12^\circ$.  Because $\angle BAC = \frac{1}{2} \cdot 72 = 36^\circ$, the ratio of the magnitudes of $\angle OBE$ and $\angle BAC$ is $\frac{12}{36} = \frac{1}{3}$, which is answer choice $\boxed{\textbf{(D)}}$.

 


	\item (\textit{AHSME 1963}) A gives $B$ as many cents as $B$ has and $C$ as many cents as $C$ has. Similarly, $B$ then gives $A$ and $C$ as many cents as each then has. 
$C$, similarly, then gives $A$ and $B$ as many cents as each then has. If each finally has $16$ cents, with how many cents does $A$ start?


 $\textbf{(A)}\ 24 \qquad \textbf{(B)}\ 26\qquad \textbf{(C)}\ 28 \qquad \textbf{(D)}\ 30 \qquad \textbf{(E)}\ 32$


 
 \textbf{Answer: }B



\textbf{Solution}

Let $a$ be number of cents $A$ originally had, $b$ be number of cents $B$ originally had, and $c$ be number of cents $C$ originally had.

 After $A$ gave his money away, $A$ has $a-b-c$ cents, $B$ has $2b$ cents, and $C$ has $2c$ cents.

 After $B$ gave his money away, $A$ has $2a-2b-2c$ cents, $B$ has $-a+3b-c$ cents, and $C$ has $4c$ cents.

 After $C$ gave his money away, $A$ has $4a-4b-4c$ cents, $B$ has $-2a+6b-2c$ cents, and $C$ has $-a-b+7c$ cents.

 Since all of them have $16$ cents in the end, we can write a \href{https://artofproblemsolving.com/wiki/index.php?title=System\_of\_equations}{system of equations}.
 \[4a-4b-4c=16\]
 \[-2a+6b-2c=16\]
 \[-a-b+7c=16\]Note that adding the three equation yields $a+b+c=48$, so $4a+4b+4c=192$.  Therefore, $8a=208$, so $a = 26$.  Solving for $a$ can also be done traditionally.

 Thus, $A$ started out with $26$ cents, which is answer choice $\boxed{\textbf{(B)}}$.

 


 



\textbf{Solution 2}

We know that people not giving away money on the previous turn now have twice what they previously had. Using the fact that the sum of their money is $48$ cents, we can work backward
 \[16,16,16\]
 \[8,8,32\]
 \[4,28,16\]
 \[26,14,8\]Thus at the beginning $A$ has $26\boxed{B}$ cents.

 


	\item (\textit{AHSME 1963}) Consider equations of the form $x^2 + bx + c = 0$. How many such equations have real roots and have coefficients $b$ and $c$ selected 
from the set of integers $\{1,2,3, 4, 5,6\}$?


 $\textbf{(A)}\ 20 \qquad \textbf{(B)}\ 19 \qquad \textbf{(C)}\ 18 \qquad \textbf{(D)}\ 17 \qquad \textbf{(E)}\ 16$


 
 \textbf{Answer: }B



\textbf{Solution}

The \href{https://artofproblemsolving.com/wiki/index.php?title=Discriminant}{discriminant} of the quadratic is $b^2 - 4c$.  Since the quadratic has real roots,
 \[b^2 - 4c \ge 0\]
 \[b^2 \ge 4c\]If $b = 6$, then $c$ can be from $1$ to $6$.  If $b = 5$, then $c$ can also be from $1$ to $6$.  If $b=4$, then $c$ can be from $1$ to $4$.  If $b=3$, then $c$ can be $1$ or $2$.  If $b=2$, then $c$ can only be $1$.  If $b = 1$, no values of $c$ in the set would work.

 Thus, there are a total of $19$ equations that work.  The answer is $\boxed{\textbf{(B)}}$.

 


	\item (\textit{AHSME 1963}) Point $F$ is taken in side $AD$ of square $ABCD$. At $C$ a perpendicular is drawn to $CF$, meeting $AB$ extended at $E$. 
The area of $ABCD$ is $256$ square inches and the area of $\triangle CEF$ is $200$ square inches. Then the number of inches in $BE$ is:


 
                    \begin{figure}[H]
                        centering
                        \includegraphics[width=0.5\textwidth]{Problem/AHSME_I_1963_25_Problem_Figure_1.png}
                    \end{figure}
                    


 $\textbf{(A)}\ 12 \qquad \textbf{(B)}\ 14 \qquad \textbf{(C)}\ 15 \qquad \textbf{(D)}\ 16 \qquad \textbf{(E)}\ 20$


 
 \textbf{Answer: }A



\textbf{Solution}

Because $ABCD$ is a \href{https://artofproblemsolving.com/wiki/index.php?title=Square}{square}, $DC = CB = 16$, $DC \perp DA$, and $CB \perp BA$.  Also, because $\angle DCF + \angle FCB = \angle FCB + \angle BCE = 90^\circ$, $\angle DCF = \angle BCE$.  Thus, by ASA Congruency, $\triangle DCF \cong \triangle BCF$.

 From the congruency, $CF = CE$.  Using the area formula for a triangle, $CE = 20$.  Finally, by the \href{https://artofproblemsolving.com/wiki/index.php?title=Pythagorean\_Theorem}{Pythagorean Theorem}, $BE = \sqrt{20^2 - 16^2} = 12$, which is answer choice $\boxed{\textbf{(A)}}$.

 


	\item (\textit{AHSME 1963}) Version 1
Consider the statements:


 $\textbf{(1)}\ p\text{ }\wedge\sim q\wedge r\qquad\textbf{(2)}\ \sim p\text{ }\wedge\sim q\wedge r\qquad\textbf{(3)}\ p\text{ }\wedge\sim q\text{ }\wedge\sim r\qquad\textbf{(4)}\ \sim p\text{ }\wedge q\text{ }\wedge r$


 where $p,q$, and $r$ are propositions. How many of these imply the truth of $(p\rightarrow q)\rightarrow r$?


 $\textbf{(A)}\ 0 \qquad \textbf{(B)}\ 1\qquad \textbf{(C)}\ 2 \qquad \textbf{(D)}\ 3 \qquad \textbf{(E)}\ 4$


 Version 2
Consider the statements (1) $p$ and $r$ are true and $q$ is false (2) $r$ is true and $p$ and $q$ are false (3) $p$ is true and $q$ and $r$ are false (4) $q$ and $r$ are true and $p$ is false. 
How many of these imply the truth of the statement
"$r$ is implied by the statement that $p$ implies $q$"?


 $\textbf{(A)}\ 0 \qquad \textbf{(B)}\ 1\qquad \textbf{(C)}\ 2 \qquad \textbf{(D)}\ 3 \qquad \textbf{(E)}\ 4$


 
 \textbf{Answer: }E



\textbf{Solution}

Statement $1$ states that $p$ is true and $q$ is false.  Therefore, $p \rightarrow q$ is false, because a premise being true and a conclusion being false is, itself, false.  This means that $(p \rightarrow q) \rightarrow X$, where $X$ is any logical statement (or series of logical statements) must be true - if your premise is false, then the implication is automatically true.  So statement $1$ implies the truth of the given statement.

 Statement $3$ similarly has $p$ as true and $q$ is false, so it also implies the truth of the given statement.

 Statement $2$ states that $p$ and $q$ are both false.  This in turn means that $p \rightarrow q$ is true.  Since $r$ is also true from statement $2$, this means that $(p \rightarrow q) \rightarrow r$ is true, since $T \rightarrow T$ is $T$.  Thus statement $2$ implies the truth of the given statement.

 Statement $4$ states that $p$ is false and $q$ is true.  In this case, $p \rightarrow q$ is true - your conclusion can be true even if your premise is false.  And, since $r$ is also true from statement $4$, this means $(p \rightarrow q) \rightarrow r$ is true.  Thus, statement $4$ implies the truth of the given statement.

 All four statements imply the truth of the given statement, so the answer is $\boxed{\textbf{(E)}}$

 


	\item (\textit{AHSME 1963}) Six straight lines are drawn in a plane with no two parallel and no three concurrent. The number of regions into which they divide the plane is:


 $\textbf{(A)}\ 16 \qquad \textbf{(B)}\ 20\qquad \textbf{(C)}\ 22 \qquad \textbf{(D)}\ 24 \qquad \textbf{(E)}\ 26$


 
 \textbf{Answer: }C



\textbf{Solution 1}

The first line divides the plane into two regions.  The second line intersects one line, creating two regions.  The third line intersects two lines, creating three regions.  Similarly, the fourth line intersects three lines and creates four regions, the fifth line intersects four lines and creates five regions, and the sixth line intersects five lines and creates six regions.

 Totaling the regions created results in $2 + 2 + 3 + 4 + 5 + 6 = 22$ regions, which is answer choice $\boxed{\textbf{(C)}}$.

 



\textbf{Solution 2}


                    \begin{figure}[H]
                        centering
                        \includegraphics[width=0.5\textwidth]{Solution/AHSME_I_1963_27_Solution_Figure_1.png}
                    \end{figure}
                    

 With careful drawing, one can draw all six lines and count the regions.  There are $22$ regions in total, which is answer choice $\boxed{\textbf{(C)}}$.

 



\textbf{Solution 3}

We can use the fact that the number of regions that $n$ lines divide a plane is given by the equation $L_n = \frac{n^2 + n +2}{2}$, and in this problems, $n=6$, from which the answer is $\boxed{22}$. 

 The formula is proved in Engel's Extremal Principles chapter or \href{https://www.cut-the-knot.org/proofs/LinesDividePlane.shtml}{https://www.cut-the-knot.org/proofs/LinesDividePlane.shtml} gives two proofs for it.

 (OmicronGamma 2019)

 


	\item (\textit{AHSME 1963}) Given the equation $3x^2 - 4x + k = 0$ with real roots. The value of $k$ for which the product of the roots of the equation is a maximum is:


 $\textbf{(A)}\ \frac{16}{9}\qquad \textbf{(B)}\ \frac{16}{3}\qquad \textbf{(C)}\ \frac{4}{9}\qquad \textbf{(D)}\ \frac{4}{3}\qquad \textbf{(E)}\ -\frac{4}{3}$


 
 \textbf{Answer: }D



\textbf{Solution}

By \href{https://artofproblemsolving.com/wiki/index.php?title=Vieta\%27s\_Formulas}{Vieta's Formulas}, the product of the roots is $\frac{k}{3}$.  This value increases as $k$ increases.

 Also, the quadratic’s roots are real, then the \href{https://artofproblemsolving.com/wiki/index.php?title=Discriminant}{discriminant} is greater than or equal to zero, so
 \[16-12k \ge 0\]
 \[-12k \ge -16\]
 \[k \le \frac{4}{3}\]

 Thus, the $k$ value that maximizes the product of the roots is $\frac{4}{3}$, which is answer choice $\boxed{\textbf{(D)}}$.

 


 


	\item (\textit{AHSME 1963}) A particle projected vertically upward reaches, at the end of $t$ seconds, an elevation of $s$ feet where $s = 160 t - 16t^2$. The highest elevation is:


 $\textbf{(A)}\ 800 \qquad \textbf{(B)}\ 640\qquad \textbf{(C)}\ 400 \qquad \textbf{(D)}\ 320 \qquad \textbf{(E)}\ 160$


 
 \textbf{Answer: }C



\textbf{Solution}

The highest elevation a particle can reach is the vertex of the quadratic.  The x-value that can get the maximum is $\frac{-160}{-2 \cdot 16} = 5$, so the highest elevation is $160(5) - 16(5^2) = 400$ feet, which is answer choice $\boxed{\textbf{(C)}}$.

 


	\item (\textit{AHSME 1963}) Let $F=\log\dfrac{1+x}{1-x}$. Find a new function $G$ by replacing each $x$ in $F$ by $\dfrac{3x+x^3}{1+3x^2}$, and simplify. 
The simplified expression $G$ is equal to:


 $\textbf{(A)}\ -F \qquad \textbf{(B)}\ F\qquad \textbf{(C)}\ 3F \qquad \textbf{(D)}\ F^3 \qquad \textbf{(E)}\ F^3-F$


 
 \textbf{Answer: }C



\textbf{Solution}

Replace the $x$ in $F$ with $\frac{3x+x^3}{1+3x^2}$.
 \[\log\dfrac{1 + \frac{3x+x^3}{1+3x^2} }{1 - \frac{3x+x^3}{1+3x^2} }\]
 \[\log\dfrac{\frac{1+3x^2}{1+3x^2} + \frac{3x+x^3}{1+3x^2} }{\frac{1+3x^2}{1+3x^2} - \frac{3x+x^3}{1+3x^2} }\]
 \[\log (\frac{1+3x+3x^2+x^3}{1+3x^2} \div \frac{1-3x+3x^2-x^3}{1+3x^2})\]Recall that from the \href{https://artofproblemsolving.com/wiki/index.php?title=Binomial\_Theorem}{Binomial Theorem}, $(1+x)^3 = 1+3x+3x^2+x^3$ and $(1-x)^3 = 1-3x+3x^2-x^3$.
 \[\log (\frac{(1+x)^3}{1+3x^2} \div \frac{(1-x)^3}{1+3x^2})\]
 \[\log\dfrac{(1+x)^3}{(1-x)^3}\]
 \[3 \log\dfrac{1+x}{1-x}\]Thus, the expression is equivalent to $3F$, which is answer choice $\boxed{\textbf{(C)}}$.

 


	\item (\textit{AHSME 1963}) The number of solutions in positive integers of $2x+3y=763$ is:


 $\textbf{(A)}\ 255 \qquad \textbf{(B)}\ 254\qquad \textbf{(C)}\ 128 \qquad \textbf{(D)}\ 127 \qquad \textbf{(E)}\ 0$


 
 \textbf{Answer: }D



\textbf{Solution 1}

Solving for $x$ in the equation yields $x =rfthe meaning of theta 0$.  Solving the inequality results in $y \le 254 \frac{1}{3}$.  From the two conditions, $y$ can be an odd number from $1$ to $253$, so there are $127$ solutions where $x$ and $y$ are integers.  The answer is $\boxed{\textbf{(D)}}$.

 



\textbf{Solution 2}

We will prove that $y$ is an odd number by contradiction. If $y$ is even, then we know that $y = 2m$ where $m$ is some integer. However, this immediately assumes that $\text{ even } + \text{ even } = \text{ odd }$ which is impossible. therefore $y$ must ben odd. then we can easily prove that $x$ .....

 



\textbf{Solution 3}

We can solve for first solution by applying extended euclids division algorithm or we can apply hit and trial for the first solution to get $x_0$ = $380$ and $y_0$ =$1$ . then the general solution of the given diophanitine equation will be $x$ = $x_0$ +$3t$ and $y$ = $y_0$ - $2t$. Since we need only positive integer solutions So we solve $380$ + $3t$ $>$ $0$ and $1$-$2t$ $>$ $0$ to get $t$ $>$ $0$ (applying Greatest integer function) also we can clearly see that $t_{(min)}$ $=$ $0$ so,t $<$ $GIF$($383$/$3$). That implies $t$ ranges from $0$ to $127$. Hence,the correct answer is $127$, $\boxed{\textbf{(D)}}$. 

 


	\item (\textit{AHSME 1963}) The dimensions of a rectangle $R$ are $a$ and $b$, $a < b$. It is required to obtain a rectangle with dimensions $x$ and $y$, $x < a, y < a$, 
so that its perimeter is one-third that of $R$, and its area is one-third that of $R$. The number of such (different) rectangles is:


 $\textbf{(A)}\ 0 \qquad \textbf{(B)}\ 1\qquad \textbf{(C)}\ 2 \qquad \textbf{(D)}\ 4 \qquad \textbf{(E)}\ \infty$


 
 \textbf{Answer: }A



\textbf{Solution}

Using the perimeter and area formulas,
 \[2(x+y) = \frac{2}{3}(a+b)\]
 \[x+y = \frac{a+b}{3}\]
 \[xy = \frac{ab}{3}\]Dividing the second equation by the last equation results in
 \[\frac1y + \frac1x = \frac1b + \frac1a\]Since $x,y < a$, $\frac1a < \frac1x, \frac1y$.  Since $a < b$, $\frac1b < \frac1a$.  That means
 \[\frac1x + \frac1y > \frac1a + \frac1a > \frac1a + \frac1b\]This is a contradiction, so there are $\boxed{\textbf{(A)}\ 0}$ rectangles that satisfy the conditions.

 


	\item (\textit{AHSME 1963}) Given the line $y = \dfrac{3}{4}x + 6$ and a line $L$ parallel to the given line and $4$ units from it. A possible equation for $L$ is:


 $\textbf{(A)}\ y =\frac{3}{4}x+1\qquad \textbf{(B)}\ y =\frac{3}{4}x\qquad \textbf{(C)}\ y =\frac{3}{4}x-\frac{2}{3}\qquad  \\ \textbf{(D)}\ y = \dfrac{3}{4}x -1 \qquad \textbf{(E)}\ y = \dfrac{3}{4}x + 2$


 
 \textbf{Answer: }A



\textbf{Solution}


                    \begin{figure}[H]
                        centering
                        \includegraphics[width=0.5\textwidth]{Solution/AHSME_I_1963_33_Solution_Figure_1.png}
                    \end{figure}
                    

 If a point on line $L$ is four units away from another line, the perpendicular distance between the two lines is $4$.  Since a line perpendicular to the original has a slope of $-\frac{4}{3}$, a point on that line is $5d$ units away and can be derived from subtracting $4d$ from the y-coordinate and adding $3d$ to the x-coordinate of a point on the original line.

 We want the perpendicular distance to be $4$, so $d = \frac{4}{5}$.  Since one point on the original line is $(0,6)$, a point on line $L$ will be $(0+3 \cdot \frac{4}{5},6 - 4 \cdot \frac{4}{5})$, which is simplified as $(\frac{12}{5},\frac{14}{5})$.  Using point-slope form, the equation of line $L$ is
 \[y - \frac{14}{5} = \frac{3}{4} (x - \frac{12}{5})\]
 \[y = \frac{3}{4} x + 1\]The answer is $\boxed{\textbf{(A)}}$.

 


	\item (\textit{AHSME 1963}) In $\triangle ABC$, side $a = \sqrt{3}$, side $b = \sqrt{3}$, and side $c > 3$. Let $x$ be the largest number such that the magnitude, 
in degrees, of the angle opposite side $c$ exceeds $x$. Then $x$ equals:


 $\textbf{(A)}\ 150^{\circ} \qquad \textbf{(B)}\ 120^{\circ}\qquad \textbf{(C)}\ 105^{\circ} \qquad \textbf{(D)}\ 90^{\circ} \qquad \textbf{(E)}\ 60^{\circ}$


 
 \textbf{Answer: }B



\textbf{Solution}

Using the \href{https://artofproblemsolving.com/wiki/index.php?title=Law\_of\_Cosines}{Law of Cosines},
 \[\sqrt{3 + 3 - 2\cdot 3 \cdot \cos{x^\circ}}>3\]

 Both sides are positive, so squaring both sides will not affect the inequality.

 
 \[6 - 6 \cos{x^\circ}>9\]
 \[\cos{x^\circ} < -\frac{1}{2}\]

 Note that $\cos{120^\circ} = -\frac{1}{2}$.  As $x$ gets closer to $180^{\circ}$, $\cos{x}$ decreases towards $-1$.  Thus, $x \leq 120$, and in order to maximize $x$, the answer has to be $\boxed{\textbf{(B)}}$

 


	\item (\textit{AHSME 1963}) The lengths of the sides of a triangle are integers, and its area is also an integer. 
One side is $21$ and the perimeter is $48$. The shortest side is:


 $\textbf{(A)}\ 8 \qquad \textbf{(B)}\ 10\qquad \textbf{(C)}\ 12 \qquad \textbf{(D)}\ 14 \qquad \textbf{(E)}\ 16$


 
 \textbf{Answer: }B



\textbf{Solution}

Let $b$ and $c$ be the other two sides of the triangle.  The \href{https://artofproblemsolving.com/wiki/index.php?title=Perimeter}{perimeter} of the triangle is $48$ units, so $c = 27-b$ and the semiperimeter equals $24$ units.

 By \href{https://artofproblemsolving.com/wiki/index.php?title=Heron\%27s\_Formula}{Heron's Formula}, the area of the triangle is $\sqrt{24 \cdot 3(24-b)(b-3)}$.  Plug in the answer choices for $b$ and write the \href{https://artofproblemsolving.com/wiki/index.php?title=Prime\_factorization}{prime factorization} of the product to make sure it is a perfect square.

 Testing $b = 8$ results in the area being $\sqrt{6 \cdot 4 \cdot 3 \cdot 16 \cdot 5} = \sqrt{2^7 \cdot 3^2 \cdot 5}$, so $8$ does not work.  However, testing $b = 10$ results in the area being $\sqrt{6 \cdot 4 \cdot 3 \cdot 14 \cdot 7} = \sqrt{2^4 \cdot 3^2 \cdot 7^2}$, so $10$ works.  The third side is $17$, and the sides satisfy the \href{https://artofproblemsolving.com/wiki/index.php?title=Triangle\_Inequality}{Triangle Inequality}, so the answer is $\boxed{\textbf{(B)}}$.

 


	\item (\textit{AHSME 1963}) A person starting with $\$64$ and making $6$ bets, wins three times and loses three times, 
the wins and losses occurring in random order. The chance for a win is equal to the chance for a loss. 
If each wager is for half the money remaining at the time of the bet, then the final result is:


 $\textbf{(A)}\text{ a loss of }\$ 27 \qquad \textbf{(B)}\text{ a gain of }\$ 27 \qquad \textbf{(C)}\text{ a loss of }\$ 37 \qquad \\ \textbf{(D)}\text{ neither a gain nor a loss}\qquad \\ \textbf{(E)}\text{ a gain or a loss depending upon the order in which the wins and losses occur}$


 
 \textbf{Answer: }C



\textbf{Solution}

If the person wins the bet, the person has $\frac{3}{2}$ of the previous amount.  If the person loses the bet, the person only has $\frac{1}{2}$ of the previous amount.

 Because of the Commutative Property, the order of multiplying the multipliers does not matter.  Thus, the person walks away with $64 \cdot \frac{3}{2} \cdot \frac{3}{2} \cdot \frac{3}{2} \cdot \frac{1}{2} \cdot \frac{1}{2} \cdot \frac{1}{2} = 27$ dollars, so the person loses $\$ 37$.  The answer is $\boxed{\textbf{(C)}}$.

 


	\item (\textit{AHSME 1963}) Given points $P_1, P_2,\cdots,P_7$ on a straight line, in the order stated (not necessarily evenly spaced). 
Let $P$ be an arbitrarily selected point on the line and let $s$ be the sum of the undirected lengths
$PP_1, PP_2, \cdots , PP_7$. Then $s$ is smallest if and only if the point $P$ is:


 $\textbf{(A)}\ \text{midway between }P_1\text{ and }P_7\qquad \\ \textbf{(B)}\ \text{midway between }P_2\text{ and }P_6\qquad \\ \textbf{(C)}\ \text{midway between }P_3\text{ and }P_5\qquad \\ \textbf{(D)}\ \text{at }P_4 \qquad \textbf{(E)}\ \text{at }P_1$


 
 \textbf{Answer: }D



\textbf{Solution}

By the \href{https://artofproblemsolving.com/wiki/index.php?title=Triangle\_Inequality}{Triangle Inequality}, $P_1P + P_7P \ge P_1P_7$, with equality happening when $P$ is between $P_1$ and $P_7$. Using similar logic, $P$ must be between $P_3$ and $P_5$ in order for the distance to be minimized.

 The only point left to deal with is $P_4$ (which is also between $P_3$ and $P_5$).  The minimum possible distance is $0$ (when $P$ is on $P_4$), so the answer is $\boxed{\textbf{(D)}}$.

 


	\item (\textit{AHSME 1963}) Point $F$ is taken on the extension of side $AD$ of parallelogram $ABCD$. $BF$ intersects diagonal $AC$ at $E$ and side $DC$ at $G$. 
If $EF = 32$ and $GF = 24$, then $BE$ equals:


 
                    \begin{figure}[H]
                        centering
                        \includegraphics[width=0.5\textwidth]{Problem/AHSME_I_1963_38_Problem_Figure_1.png}
                    \end{figure}
                    


 $\textbf{(A)}\ 4 \qquad \textbf{(B)}\ 8\qquad \textbf{(C)}\ 10 \qquad \textbf{(D)}\ 12 \qquad \textbf{(E)}\ 16$


 
 \textbf{Answer: }E



\textbf{Solution}

Let $BE = x$ and $BC = y$.  Since $AF \parallel BC$, by AA Similarity, $\triangle AFE \sim \triangle CBE$.  That means $\frac{AF}{CB} = \frac{FE}{BE}$.  Substituting in values results in
 \[\frac{AF}{y} = \frac{32}{x}\]Thus, $AF = \frac{32y}{x}$, so $FD = \frac{32y - xy}{x}$.

 
In addition, $DC \parallel AB$, so by AA Similarity, $\triangle FDG = \triangle FAB$.  That means 
 \[\frac{\frac{32y-xy}{x}}{\frac{32y}{x}} = \frac{24}{x+32}\]Cross multiply to get
 \[\frac{y(32-x)}{x} (x+32) = \frac{32y}{x} \cdot 24\]Since $x \ne 0$ and $y \ne 0$,
 \[(32-x)(32+x) = 32 \cdot 24\]
 \[32^2 - x^2 = 768\]
 \[1024 - 768 = x^2\]
 \[256 = x^2\]
 \[16,-16 = x^2\]As the lenght of lines can never be negative, $x = 16$, which is answer choice $\boxed{\textbf{(E)}}$.

 


	\item (\textit{AHSME 1963}) In $\triangle ABC$ lines $CE$ and $AD$ are drawn so that 
$\dfrac{CD}{DB}=\dfrac{3}{1}$ and $\dfrac{AE}{EB}=\dfrac{3}{2}$. Let $r=\dfrac{CP}{PE}$ 
where $P$ is the intersection point of $CE$ and $AD$. Then $r$ equals:


 
                    \begin{figure}[H]
                        centering
                        \includegraphics[width=0.5\textwidth]{Problem/AHSME_I_1963_39_Problem_Figure_1.png}
                    \end{figure}
                    


 

$\textbf{(A)}\ 3 \qquad \textbf{(B)}\ \dfrac{3}{2}\qquad \textbf{(C)}\ 4 \qquad \textbf{(D)}\ 5 \qquad \textbf{(E)}\ \dfrac{5}{2}$


 
 \textbf{Answer: }D



\textbf{Solution 1}


                    \begin{figure}[H]
                        centering
                        \includegraphics[width=0.5\textwidth]{Solution/AHSME_I_1963_39_Solution_Figure_1.png}
                    \end{figure}
                    

 Draw line $PB$, and let $[PEB] = 2b$, $[PDB] = a$, and $[CAP] = c$, so $[CPD] = 3a$ and $[APE] = 3b$.  Because $\triangle CAE$ and $\triangle CEB$ share an altitude,
 \[c + 3b = \frac{3}{2} (3a+a+2b)\]
 \[c + 3b = 6a + 3b\]
 \[c  = 6a\]Because $\triangle ACD$ and $\triangle ABD$ share an altitude,
 \[6a+3a = 3(a+2b+3b)\]
 \[9a = 3a+15b\]
 \[6a = 15b\]
 \[a = \frac{5}{2}b\]Thus, $[CAP] = 15b$, and since $[APE] = 3b$, $r = \frac{CP}{PE} = 5$, which is answer choice $\boxed{\textbf{(D)}}$.

 



\textbf{Solution 2 (Mass Geometry)}

Let the mass of point $A$, $B$, $C$, $D$, and $E$ be $mA$, $mB$, $mC$, $mD$, and $mE$ respectively.By mass geometry theorems, we have
 \begin{align*} \frac{CP}{PE} = \frac{mE}{mC} \end{align*}
Focusing on the line segment $AB$, using mass geometry theorems, we have
 \begin{align*} mE = mA + mB \end{align*}
and
 \begin{align*} \frac{3}{2} &= \frac{mB}{mA} \\ mA &= \frac{2}{3}mB \end{align*}
which leads to $mE = \frac{5}{3}mB$.For line segment $CB$, similarly, we got
 \[mC = \frac{1}{3}mB\]Substituting $mC$ and $mE$ back to the equation we obtained at the beginning, we got:
 \begin{align*} \frac{CP}{PE} = \frac{mE}{mC} = \frac{\frac{5}{3}mB}{\frac{1}{3}mB} = 5 \end{align*}
which gives us the answer choice $\boxed{\textbf{(D)}}$. -nullptr07

 



\textbf{Solution 3 (Menelaus’s Theorem)}

By using Menelaus’s Theorem on triangle BCE and points A, P, D, we can substitute in the known values to find that the answer is 5.-pengu14

 



\textbf{Solution 4 (Ceva and Menelaus)}

Draw $BP$ to intersect line $CA$ at $X$. Then, through Ceva's on $ABC$,

 
 \[\frac{CX}{XA} \cdot \frac{AE}{EB} \cdot \frac{BD}{DC} = 1\]
 \[\frac{CX}{XA} \cdot \frac{3}{2} \cdot \frac{1}{3} = 1\]
 \[\frac{CX}{XA} = 2\].

 Then, through Menelaus' on $CAE$,

 
 \[\frac{CX}{XA} \cdot \frac{AB}{BE} \cdot \frac{EP}{PC} = 1\]
 \[2 \cdot \frac{5}{2} \cdot \frac{EP}{PC} = 1\]
 \[\frac{EP}{PC} = \frac{1}{5}\]

 Thus $\frac{PC}{EP} = \boxed{5}$

 


	\item (\textit{AHSME 1963}) If $x$ is a number satisfying the equation $\sqrt[3]{x+9}-\sqrt[3]{x-9}=3$, then $x^2$ is between:


 $\textbf{(A)}\ 55\text{ and }65\qquad \textbf{(B)}\ 65\text{ and }75\qquad \textbf{(C)}\ 75\text{ and }85\qquad \textbf{(D)}\ 85\text{ and }95\qquad \textbf{(E)}\ 95\text{ and }105$


 
 \textbf{Answer: }C



\textbf{Solution 1}

Let $a = \sqrt[3]{x + 9}$ and $b = \sqrt[3]{x - 9}$. Cubing these equations, we get $a^3 = x + 9$ and $b^3 = x - 9$, so $a^3 - b^3 = 18$. The left-hand side factors as
 \[(a - b)(a^2 + ab + b^2) = 18.\]

 However, from the given equation $\sqrt[3]{x + 9} - \sqrt[3]{x - 9} = 3$, we get $a - b = 3$. Then $3(a^2 + ab + b^2) = 18$, so $a^2 + ab + b^2 = 18/3 = 6$.

 Squaring the equation $a - b = 3$, we get $a^2 - 2ab + b^2 = 9$. Subtracting this equation from the equation $a^2 + ab + b^2 = 6$, we get $3ab = -3$, so $ab = -1$. But $a = \sqrt[3]{x + 9}$ and $b = \sqrt[3]{x - 9}$, so $ab = \sqrt[3]{(x + 9)(x - 9)} = \sqrt[3]{x^2 - 81}$, so $\sqrt[3]{x^2 - 81} = -1$. Cubing both sides, we get $x^2 - 81 = -1$, so $x^2 = 80$. The answer is $\boxed{\textbf{(C)}}$.

 


 



\textbf{Solution 2}

$\sqrt[3]{x+9}-\sqrt[3]{x-9}-3=0$ i.e, $\sqrt[3]{x+9}+\sqrt[3]{-x+9}+(-3)=0$

 \begin{verbatim}
if the sum of three numbers is zero, then their sum of cubes is thrice the product of each number.then, $x+9-x+9+27=3(\sqrt[3]{x+9})(\sqrt[3]{9-x})(-3)$ . by solving this, we get $-9=-9\sqrt[3]{9^2-x^2}$$x^2=80$. this gives the step what we had done in solution 1.The answer is $\boxed{\textbf{(C)}}$.
\end{verbatim}





\textbf{Solution 3}

Cubing both sides, we get
 \[x+9-3\sqrt[3]{(x+9)^2(x-9)}+3\sqrt[3]{(x+9)(x-9)^2}-x+9=27,\]so 
 \[18-3\sqrt[3]{(x-9)(x+9)}(\sqrt[3]{x+9}-\sqrt[3]{x-9})=27\implies 3\sqrt[3]{x^2-81}=-3.\] Dividing both sides by 3 and cubing, we find $x^2=80$, which is between $\boxed{(C)75\text{ and }85}$.

 



\textbf{Solution 4}

We consider the formula $(a-b)^3 = a^3 - 3a^2b + 3ab^2 - b^3$. Factoring out $3ab$, and the commutative property gives $(   a^3 - b^3) - 3ab(a-b)$. Cubing both sides gives us
 \[(\sqrt[3]{x+9}-\sqrt[3]{x-9})^3=27\]Using our formula, we have
 \[18 - 9 \sqrt[3]{x^2 - 81} = 27\]Solving this gives $x^2 = 80$, therefore, the answer is $\boxed{\textbf{(C)}}$.~zixuan12

 


\end{enumerate}